\documentclass[11pt]{article}
% Standalone preamble for web-only supplementary notes.
\usepackage[T1]{fontenc}
\usepackage{amsmath,amssymb,amsthm}
\usepackage{hyperref}
\usepackage{cleveref}
\DeclareMathOperator{\Map}{Map}
\DeclareMathOperator{\Fun}{Fun}
\newcommand{\iso}{\xrightarrow{\sim}}
\newtheorem{lemma}{Supplementary Lemma}
\newtheorem{corollary}[lemma]{Supplementary Corollary}
\crefname{lemma}{Supplementary Lemma}{Supplementary Lemmas}
\Crefname{lemma}{Supplementary Lemma}{Supplementary Lemmas}
\crefname{corollary}{Supplementary Corollary}{Supplementary Corollaries}
\Crefname{corollary}{Supplementary Corollary}{Supplementary Corollaries}

\begin{document}
\section*{Reordering inverse conditions}

In the \href{rezk-axiom.html}{construction of inverse data}, a triangle
has one short edge prescribed and its long edge identified with an identity.
One can impose these constraints in either order. The comparison must
preserve the projection to the original family of arrows.

The first two lemmas use only the absolute pullback structure of Chapter 1.

\begin{lemma}
\label[lemma]{lem:Supplementary_Two_Stage_Fibers}
Let $f\colon B\to A$, $g\colon T\to A$, $h\colon T\to E$,
and $s\colon D\to E$. Put $R=B\times_A T$, with projections $p,q$,
and $J=T\times_E D$, with projection $j\colon J\to T$.
There is a specified equivalence
\[
 \Phi\colon R\times_E D\longrightarrow J\times_A B,
 \qquad \pi_B\Phi=p\pi_R,
\]
where the two fiber products use $(hq,s)$ and $(gj,f)$ respectively.
In particular, if $p\pi_R$ is an equivalence, so is $\pi_B$.
\end{lemma}
\begin{proof}[Proof idea]
Nested pullback comparison first gives
$R\times_E D\simeq R\times_T J$. Exchange the two legs to obtain
$J\times_T R$, then flatten using the defining pullback of $R$, with
its legs exchanged, to obtain $J\times_A B$.
All three comparison functors are equivalences by pullback pasting.

Trace the projection to $B$ through the three stages. Their specified
projection comparisons give $\pi_B\Phi=p\pi_R$.
The two-out-of-three property then gives the last assertion. This step
uses the projection comparison in addition to equivalence of the categories.
\end{proof}

\begin{lemma}
\label[lemma]{lem:Supplementary_Section_Fiber}
Let $e\colon E\to C$ and $c\colon C\to E$ be equipped with an
identification $ec=\mathrm{id}_C$. Let $T$ be a pullback of $e$ along
$b\colon B\to C$, with projections $h\colon T\to E$ and
$t\colon T\to B$. Then
\[
 t\pi_T\colon T\times_E C\longrightarrow B
\]
 is an equivalence, where the pullback uses $h$ and $c$.
\end{lemma}
\begin{proof}[Proof idea]
Paste the defining pullback of $T\times_E C$ with the given pullback
square for $T$. The outer rectangle is a pullback of $ec$ along $b$.
Since $ec$ is identified with the identity, its pullback projection is
an equivalence. Pullback symmetry supplies the orientation used in the
nested construction.
\end{proof}

\section*{Application to inverse triangles}

Write $\operatorname{Triang}(C)$ for the triangle category,
$\operatorname{Ar}(C)$ for the arrow category, and $d_C\colon C\to
\operatorname{Ar}(C)$ for the identity-arrow functor. Put
\[
 J(C)=\operatorname{Triang}(C)\times_{\operatorname{Ar}(C)} C,
\]
using long-edge restriction and $d_C$. Its two short-edge maps to
$\operatorname{Ar}(C)$ describe the given arrow in one-sided inverse data.

\begin{lemma}
\label[lemma]{lem:Supplementary_Inverse_Constraint_Projection}
Assume the Segal and commutative-square axioms and their interval endpoint
data. Let $f\colon B\to\operatorname{Ar}(C)$ be a family. Its canonical
morphism expression has endpoints $e_0f,e_1f$ and identity frames, where
$e_0,e_1$ are endpoint evaluations. Suppose this expression is equipped
with invertibility data. For either
short-edge map $J(C)\to\operatorname{Ar}(C)$, the projection
\[
 J(C)\times_{\operatorname{Ar}(C)}B\longrightarrow B
\]
 is an equivalence.
\end{lemma}
\begin{proof}[Proof idea]
Fix the relevant short edge first. Composition by that invertible arrow
is an equivalence, so the corresponding long-edge square is a pullback.
Apply \Cref{lem:Supplementary_Section_Fiber} using the identity-arrow
section of the relevant endpoint evaluation functor. This proves the
projection is an equivalence when the long edge is made constant second.

Apply \Cref{lem:Supplementary_Two_Stage_Fibers} to reverse the order of
the two constraints. Its projection equation identifies the tested map
with the displayed projection from one-sided inverse data. This proves
the assertion for each short edge. The Rezk axiom is not used here.
\end{proof}

The later isomorphism-embedding argument combines the two inverse
conditions. This note isolates the pullback calculation supplying their
individual projections.
\end{document}
